\subsection{Python}\label{sec:python}
\subsubsection{命名規約}
\begin{table}[H]
  \centering
  \caption{Pythonの命名規約}
  \label{tab:code-python}
  \begin{tabular}{|c|c|c|} \hline
    対象     & ケース                & 例                \\ \hline \hline
    モジュール名 & snake\_case        & database\_models \\ \hline
    パッケージ名 & snake\_case        & my\_package      \\ \hline
    クラス名   & UpperCamelCase     & UserCreate       \\ \hline
    関数名    & snake\_case        & create\_user     \\ \hline
    変数名    & snake\_case        & is\_active       \\ \hline
    定数名    & UPPER\_SNAKE\_CASE & THIS\_IS\_STATIC \\ \hline
  \end{tabular}
\end{table}
また、補足で"\_"で始まる変数は、プライベート変数を表し、外部からアクセスしないようにする。
例外クラスは、"Error"で終わるようにする。
インスタンスメソッドの第一引数は、"self"とする。
クラスメソッドの第一引数は、"cls"とする。
\subsubsection{コーディングスタイル}
コーディングスタイルは~\href{https://pep8-ja.readthedocs.io/ja/latest}{\textbf{PEP8}}~に従う。
また、意識せずともコーディングスタイルを統一するために、black Formatter,
isort, flake8を使用する。これらを保存時に自動で実行するように設定することで、
空行や改行、importの順番などを自動で整形する。
